% Generated draft: requires mathematical and source review.
\documentclass[11pt]{article}
\usepackage[margin=25mm]{geometry}
\usepackage{fontspec}
\setmainfont{Latin Modern Roman}
\usepackage{amsmath,amssymb,mathtools,unicode-math}
\setmathfont{Latin Modern Math}
\usepackage{enumitem,needspace,xurl,hyperref}
\setlength{\parindent}{0pt}
\setlength{\parskip}{5pt}
\setlength{\emergencystretch}{3em}
\newcommand{\N}{\mathbb{N}}
\newcommand{\Z}{\mathbb{Z}}
\newcommand{\Q}{\mathbb{Q}}
\newcommand{\R}{\mathbb{R}}
\newcommand{\C}{\mathbb{C}}
\newcommand{\F}{\mathbb{F}}
\newcommand{\A}{\mathbb{A}}
\newcommand{\PP}{\mathbb{P}}
\newcommand{\E}{\mathbb{E}}
\newcommand{\eps}{\varepsilon}
\newcommand{\dd}{\,\mathrm{d}}
\begin{document}
\section*{UnsolvedMath ID: 701015}
\textbf{Source problem code:} catalogue entry
\subsection{Definitions and mathematical statement}
Work with finite groups. A group is soluble if its derived series terminates at the trivial subgroup. A Fitting class is an isomorphism-closed class of groups closed under taking normal subgroups and under normal products: if \(A,B\lhd H\) and \(A,B\in\mathfrak F\), then \(AB\in\mathfrak F\). Call \(\mathfrak F\) soluble if every member is soluble.

Retain the source's technical term \emph{local Fitting class}. For a set of groups \(\mathfrak X\), let \(\mathrm{lFit}(\mathfrak X)\) denote the smallest local Fitting class containing \(\mathfrak X\), and abbreviate \(\mathrm{lFit}(\{H\})\) to \(\mathrm{lFit}(H)\). The supplied record does not define which Fitting classes are local; this technical definition, including its ambient universe, requires source clarification before the formulation is fully self-contained.

Write \(M\lhd\lhd G\) when \(M\) is subnormal in \(G\), meaning that there is a finite chain
\[
M=M_0\lhd M_1\lhd\cdots\lhd M_r=G.
\]
For each finite group \(G\), the source defines \(\Psi(G)\) to be the least normal subgroup \(N\lhd G\), with respect to inclusion, satisfying
\[
\mathrm{lFit}(N\cap M)=\mathrm{lFit}(M)
\qquad\text{for every }M\lhd\lhd G.
\]
Existence of this least subgroup is part of the supplied definition, attributed to Doerk and Hauck, and is not established here.

In the terminology of this problem, a Fitting class \(\mathfrak F\) is \emph{saturated} if
\[
\text{for every finite group }G,\qquad
\Psi(G)\in\mathfrak F\ \Longrightarrow\ G\in\mathfrak F.
\]
\textbf{Question.} Is every nonempty soluble Fitting class that is saturated in this sense a local Fitting class?

The question retains the source's explicit quantification over finite groups in the definition of saturation. Its interpretation remains conditional on clarification of the technical notion of locality.
\subsection{Short English statement}
Using the source's local Fitting closure and its associated subgroup \(\Psi(G)\), call a Fitting class saturated when membership of \(\Psi(G)\) forces membership of \(G\) for every finite group \(G\). Must every nonempty saturated Fitting class consisting of soluble groups be local? The supplied record leaves the technical definition of locality unspecified.
\subsection{Sources}
\begin{itemize}
\item Local UnsolvedMath registry, permanent ID \(701015\), exact source code \texttt{catalogue entry}, title \emph{Kourovka Notebook Problem 14.30}; supplied upstream revision \texttt{d371ee917ac9647fbb06cd6fc529deec6f3dda82}. The registry labels the problem open; this is supplied status information, not an independently verified mathematical or literature assessment.
\item \emph{The Kourovka Notebook (Unsolved Problems in Group Theory)}, original Problem 14.30, 14th issue (1999), proposed by N.~T.~Vorob'\"{e}v. The supplied background also identifies the 21st issue edition (2026), arXiv:\texttt{1401.0300v48}. The collection description names E.~I.~Khukhro and V.~D.~Mazurov as editors.
\item K.~Doerk and P.~Hauck, \emph{Arch. Math.}, volume 35, no.~3 (1980), pp.~218--227, cited in the supplied statement for the construction of \(\Psi(G)\). No article title or URL is supplied. The cited material has not been independently consulted for this draft.
\end{itemize}
\end{document}
